\chapter{Emergence and Framework Formation}
\label{ch:emergence}

\section{Why Emergence Matters}

The earlier manuscript spoke as if a framework were already present and merely
needed to be resisted or redirected. That is incomplete. Durable frameworks are
usually formed through repeated exposure, reward history, environmental
coupling, and interference from prior data. If the book does not explain how
frameworks form, it can describe failure but not origin.

This is primarily a \textbf{medium-evidence} chapter. The key modern sources
indexed for it are \emph{Metastability demystified: the foundational past, the
pragmatic present and the promising future} (Nature Reviews Neuroscience, 2024)
and \emph{Associative conditioning in gene regulatory network models increases
integrative causal emergence} (Communications Biology, 2025). These sources are
not direct measurements of study routines or personal habits. Their value is
mechanism-building: they clarify how durable higher-level organization can
arise from repeated local interaction without requiring a central controller.
Any direct transfer from these models to everyday human self-management remains
interpretive rather than strong empirical proof.

\section{Frameworks As Historical Products}

\begin{definition}[Framework formation]
\textbf{Framework formation} is the process by which repeated inputs,
behavioural rehearsals, and environmental regularities generate a relatively
stable behavioural or cognitive regime.
\end{definition}

\begin{proposition}[History-sensitive inertia]
The inertia of a current state depends not only on the present stimulus, but
also on the historical accumulation of reinforcement, expectation, and
successful reuse.
\end{proposition}

\begin{discussion}
This proposition matters because it explains why a well-trained framework often
feels ``natural.'' It is not natural in the sense of being given once and for
all. It is natural in the sense of being repeatedly rehearsed until the system
can enter it cheaply.
\end{discussion}

\section{From Local Interaction To Global Pattern}

Emergence supplies the missing layer between micro-action and macro-pattern.
Repeated low-level events can produce a higher-level structure that later feeds
back on local choice. In the language of this book:
\begin{enumerate}[nosep]
  \item repeated actions build a frame,
  \item the frame changes what later actions cost,
  \item the new cost landscape stabilizes future behaviour.
\end{enumerate}

\begin{example}[Study identity as an emergent regime]
A student who repeatedly begins work at the same place, with the same entry
cue, and the same notebook arrangement is not merely repeating tasks. The
student is building a stable regime in which the study state becomes easier to
re-enter. The framework did not pre-exist. It emerged from repetition plus
environmental consistency.
\end{example}

The older emergence literature indexed in the project, especially John
Holland's \emph{Emergence} and \emph{Introduction to the Theory of Complex
Systems}, supplies the textbook-level logic behind this example. Local rules do
not stay local. Once interaction repeats long enough, macro-organization feeds
back on the next round of micro-decisions. The business implication is simple:
if an organization or a person keeps producing the same pattern, the right
question is not only ``who chose this?'' but also ``what local rule and what
history keep regenerating it?''

\section{Metastability Rather Than Either Rigidity Or Chaos}

The Nature Reviews Neuroscience article on metastability is important because
it blocks a common mistake. Frameworks should not be imagined as either fully
rigid or fully fluid. The medium-strength claim supported by the indexed review
is that systems can sustain organized patterns while remaining able to
reconfigure. That is exactly the condition this manuscript needs. A useful
framework must be durable enough to reduce re-entry cost, but not so rigid that
it traps the agent when the task, context, or incentive changes.

This is why metastability matters more than decorative complexity language.
It gives the chapter a disciplined way to say that a regime can be stable,
fragile, and reconfigurable at the same time. For behavioural interpretation,
that explains why some routines are productive while others become cages.

\section{Associative Conditioning And Causal Emergence}

The 2025 Communications Biology paper on associative conditioning in gene
regulatory network models is used here for one medium-strength point. It shows
an explicit route by which repeated local conditioning can increase integrative
causal emergence. The manuscript does not need to copy that model into human
psychology one-for-one. It only needs the mechanism lesson: repeated local
couplings can produce a higher-level organization that later has explanatory
value in its own right.

That matters for framework formation. A framework is not merely a bag of
separate habits. Once enough local cues, reward expectations, and action
sequences become mutually reinforcing, the higher-level regime starts to govern
what feels easy, salient, and likely. In management terms, the pattern becomes
a real operating layer.

\section{Interference, Data, and Path Dependence}

The user requirement about ``past data interference'' can be made precise
enough for textbook use. The system carries traces of prior success, prior
failure, prior reward, and prior context. Those traces create path dependence:
the same present cue can produce different actions depending on what the system
has already learned to expect.

\begin{remark}
This is the point at which emergence protects the manuscript from becoming too
individualistic. Frameworks are shaped by prior data, social environment,
learned routines, and repeated contact with particular reward structures.
\end{remark}

The evidence separation matters here. Path dependence as a general systems idea
is well supported by the chapter's medium-evidence sources and textbooks.
Applying that idea to personal framework inertia is also reasonable. But any
claim that a specific human routine is caused by the exact same mechanism as a
gene regulatory model would be \textbf{speculative}. This chapter therefore
keeps the analogy at the right level of abstraction: repeated coupling creates
history-sensitive organization, and that organization changes later choice.

\section{Why Formation Matters For Intervention}

Once framework formation is understood as an emergent historical process, the
intervention logic of the book improves. Trying to defeat a bad regime only at
the moment of failure is expensive because the regime has already been built by
many earlier loops. A cheaper strategy is to redesign the repeated local
conditions from which the regime forms:
\begin{enumerate}[nosep]
  \item stabilize the desired entry cue;
  \item reduce ambiguity in the first action;
  \item make reward arrive early enough to reinforce reuse;
  \item remove repeated environmental couplings that keep rebuilding the old
        regime.
\end{enumerate}

This is the business value of emergence language. It turns the book from a
theory of resistance into a theory of production: frameworks are made, and what
is made can be remade by changing the loop that keeps producing it.

\section{Recommended Reading}

\begin{readingbox}
\textbf{Foundation}
\begin{itemize}[nosep]
  \item John Holland, \emph{Emergence}.
  \item \emph{Introduction to the Theory of Complex Systems}.
\end{itemize}
\textbf{Modern Research}
\begin{itemize}[nosep]
  \item \emph{Metastability demystified: the foundational past, the pragmatic
        present and the promising future} (Nature Reviews Neuroscience, 2024,
        medium).
  \item \emph{Associative conditioning in gene regulatory network models
        increases integrative causal emergence} (Communications Biology, 2025,
        medium).
\end{itemize}
\textbf{Frontier}
\begin{itemize}[nosep]
  \item Frontier emergence proposals should be used only after the reader
        understands historical reinforcement and path dependence.
\end{itemize}
\end{readingbox}

\begin{summarybox}
This chapter explains where frameworks come from. Its source-backed value is
mechanism-building rather than direct measurement: repeated local interactions
can generate higher-level organization, metastability explains why useful
frameworks are durable but revisable, and causal-emergence language clarifies
why a mature regime can become a real explanatory layer. That converts the
manuscript from a static model of resistance into a developmental model of
formation, consolidation, and path dependence.
\end{summarybox}
